\documentclass[twocolumn,twocolappendix,linenumbers]{aastex702}

\usepackage{amsmath}
\usepackage{bbm,bm}
\usepackage{multirow}
\usepackage{colortbl}
\usepackage{../shared/preamble}
\usepackage{hyperref}

\renewcommand{\sectionautorefname}{Section}
\renewcommand{\subsectionautorefname}{Section}
\renewcommand{\subsubsectionautorefname}{Section}
\renewcommand{\figureautorefname}{Figure}
\renewcommand{\tableautorefname}{Table}
\renewcommand{\equationautorefname}{Equation}
\renewcommand{\appendixautorefname}{Appendix}

\renewcommand{\log}{{\rm log_{\small 10}}}

\graphicspath{{./}{paper-figures/}{paper-tables/}{../../../../figures/phangs/}}

\shorttitle{TIGRESS-PHANGS Simulation Suite}
\shortauthors{Kim et al.}

\begin{document}

\title{A \PHANGS-Informed Simulation Suite for the Star-Forming Interstellar Medium:
       Sampling Design and \TIGRESS\ Initial Conditions}

\author[0000-0003-2896-3725]{Chang-Goo Kim}
\affiliation{Department of Astrophysical Sciences, Princeton University,
             4 Ivy Lane, Princeton, NJ 08544, USA}
\email{cgkim@astro.princeton.edu}

% Additional authors to be added

\begin{abstract}
We present the design of a \PHANGS-informed simulation suite for the star-forming
interstellar medium (ISM), built on the \TIGRESS\ framework. Following the philosophy
of the \CAMELS\ project in cosmology, we construct an observationally anchored
simulation ensemble that spans the range of galactic environments found in nearby
galaxies, enabling systematic comparison with observations and simulation-based
parameter inference. The \PHANGS\ megatable survey provides aperture-level measurements
of gas surface density, stellar surface density, orbital frequency, stellar scale height,
and shear across $\sim$80 nearby galaxies. We use these as an informed prior over the
five environmental parameters that enter \TIGRESS-NCR\ as initial conditions.
The sampling method fits a kernel density estimator to the observed joint distribution
of these five design fields and maps a scrambled Sobol quasi-random sequence through
the marginal quantile functions to generate the simulation design. The Sobol approach
enables progressive extension of the suite at any later stage without displacing earlier
simulations. Molecular fraction and star formation rate tracers are excluded from the
input design and reserved as validation quantities to be compared against emergent
simulation outputs.
[PLACEHOLDER: summary of simulation results and PHANGS validation to be added]
\end{abstract}

\keywords{Interstellar medium (847); Star formation (1569);
          Stellar feedback (1602); Magnetohydrodynamical simulations (1966);
          Methods: statistical (1361)}

\section{Introduction}
\label{sec:intro}

Star formation in galaxies is a multi-scale, multi-physics process whose
regulation has emerged as one of the central problems of galaxy evolution.
Across a wide range of galactic environments, observations reveal tight
correlations between local gas properties, dynamical equilibrium pressure,
and the star formation rate (SFR) surface density
\citep{2010ApJ...721..975O,2011ApJ...731...41O,2020ApJ...892..148S,2022ApJ...936..137O}.
These correlations are now understood as a signature of self-regulation:
energy and momentum injected by massive-star feedback---supernovae, stellar
winds, and radiation pressure---maintain the interstellar medium (ISM) in
an approximately steady state in which the vertical weight of the gas is
supported against gravity by turbulent, thermal, and magnetic pressures.
The resulting balance ties the local SFR to the local dynamical state of
the disk, and reproduces the observed scaling between $\Ssfr$ and the
dynamical equilibrium pressure $\PDE$ across more than four orders of
magnitude in environment
\citep{2020ApJ...892..148S,2023ApJ...945L..19S,2025ApJ...989...66V}.

The pressure-regulated, feedback-modulated (\PRFM) framework
\citep{2022ApJ...936..137O} provides a theoretical foundation for this
self-regulation. In its compact form, vertical equilibrium fixes the
midplane total pressure to the weight of the overlying gas, $\Ptot = \W$,
where $\W$ is the integrated weight of the gas column in the combined
gravitational potential of the gas, stars, and dark matter. The feedback
yield $\Ytot \equiv \Ptot/\Ssfr$ then encapsulates how efficiently young
stellar populations convert star formation into the pressure required to
support the disk. Equating these two relations yields the prediction
\begin{equation}
  \Ssfr = \frac{\W}{\Ytot},
  \label{eq:prfm}
\end{equation}
in which $\W$ is dictated by the local galactic environment and $\Ytot$ is
set by feedback physics. The \TIGRESS\ simulation framework
\citep{2024ApJ...972...67K} resolves the multiphase ISM in local
shearing-box magnetohydrodynamic (MHD) simulations with explicit
non-equilibrium cooling/heating, UV radiation transfer, cosmic-ray
transport, and clustered supernova feedback, and delivers calibrated values
of $\Ytot$ and its thermal, turbulent, and magnetic components across
metallicities. \PRFM\ has been adopted as the star formation prescription
in cosmological-volume simulations \citep{2024ApJ...975..151H} and
isolated-galaxy simulations \citep{2024ApJ...975..113J}, where it provides
a physically motivated alternative to empirical Kennicutt--Schmidt-type
prescriptions.

The observational counterpart to \TIGRESS\ is provided by the Physics at
High Angular resolution in Nearby GalaxieS (\PHANGS) survey
\citep{2021ApJS..255...19L,2023AJ....166..145L}, which delivers
spatially resolved, multi-wavelength measurements of $\sim80$ nearby
star-forming disk galaxies. \PHANGS\ aperture catalogs combine high-resolution
CO mapping with matched \HI, stellar mass, and SFR tracers to produce
joint constraints on gas, stars, and star formation at sub-kpc scales
\citep{2022AJ....164...43S,2023ApJ...945L..19S}.
Together with stellar dynamical modeling of the same galaxies
\citep{2020ApJ...892..148S,2025ApJ...989...66V}, the survey defines an
empirical envelope of local galactic environments---spanning more than two
orders of magnitude in gas and stellar surface density, and the full range
of orbital frequencies found in disk-dominated systems---against which any
ISM theory must be tested.

To date, however, comparisons between \TIGRESS-type simulations and
\PHANGS-type observations have relied almost exclusively on a handful of
representative simulations selected by hand at the solar neighborhood and a
few extreme environments. No systematic survey has yet been carried out
that spans the full \PHANGS\ parameter volume with controlled, calibrated
simulations of the resolved multiphase ISM. An analogous gap in
cosmology---between hand-picked hydrodynamic simulations and the
high-dimensional space of cosmological and astrophysical parameters
required for inference---was closed by the Cosmology and Astrophysics with
MachinE Learning Simulations (\CAMELS) project
\citep{2021ApJ...915...71V}. \CAMELS\ demonstrated that a structured
simulation ensemble covering parameter space in a controlled way enables
emulator construction, simulation-based inference, and systematic
comparison with observations, even at the modest individual-simulation
fidelity that suite-scale projects can afford. The same logic motivates a
\CAMELS-style ensemble for the star-forming ISM.

In this paper, the first in a series, we present the design of a
\PHANGS-informed simulation suite for the star-forming ISM based on the
\TIGRESS\ framework. We use the \PHANGS\ megatable
\citep{2022AJ....164...43S} to construct an observational prior over the
five environmental parameters that enter \TIGRESS\ as initial conditions:
the gas surface density $\Sgas$, stellar surface density $\Sstar$,
stellar scale height $\Hstar$, orbital frequency $\Omrot$, and shear
parameter $\qshear$. We construct the simulation design by fitting a
kernel density estimator (KDE) to the joint \PHANGS\ distribution of these
five fields and mapping a scrambled Sobol quasi-random sequence through
its marginal quantile functions. The Sobol construction allows the suite
to be extended progressively---doubling the sample size without displacing
or re-running earlier simulations---and yields low-discrepancy coverage of
the \PHANGS\ environmental volume at every intermediate sample size. The
molecular fraction $\fmol$ and SFR surface density $\Ssfr$ are
deliberately excluded from the input design and reserved as validation
quantities to be compared against emergent simulation outputs. In this
first suite, all physics parameters (metallicity, dust-to-gas ratio,
feedback prescription) are held at fiducial solar-neighborhood values;
extension to physics-parameter variations and the use of this suite for
simulation-based inference (\SBI) are the subjects of companion papers
\CGK{(Paper~II, in prep.)}.

The remainder of this paper is organized as follows. \autoref{sec:data}
describes the \PHANGS\ megatable data, the aperture-level join, and the
derivation of the five design fields and validation fields used throughout.
\CGK{Section~3} presents the sampling design, including the KDE fit, the
Sobol construction, and the fairness diagnostics used to assess marginal
and joint coverage. \autoref{sec:tigress_mapping} describes the mapping
from \PHANGS\ design fields to \TIGRESS\ initial conditions and the
fiducial physics choices for the first suite. \autoref{sec:results}
presents the resulting simulation design, an initial set of \TIGRESS\ runs,
and validation against \PHANGS-observed $\fmol$, $\Ssfr$, $\PDE$, and
$\seff$. \CGK{Section~6} discusses extensions to physics-parameter
variations, the role of this suite in enabling \SBI\ on the resolved ISM,
and connections to large-volume cosmological simulations adopting \PRFM.

\section{PHANGS Data and Derived Quantities}
\label{sec:data}

This section describes the observational dataset that anchors the
simulation design. \autoref{sec:phangs_megatable} summarizes the \PHANGS\
megatable survey and the aperture-level join used to construct a single
matched table of gas, stellar, and SFR properties. \autoref{sec:derived}
describes the derivation of the five \PRFM\ design fields from the
\PHANGS\ aperture measurements and stellar/dynamical models.
\autoref{sec:quality_cuts} states the quality cuts applied to define the
reference set $\Dref$ that drives the sampling design in
\CGK{Section~3}. \autoref{tab:fields} summarizes the design and validation
fields with their units and roles.

\subsection{PHANGS Megatable Survey}
\label{sec:phangs_megatable}

We use the publicly released \PHANGS\ megatable products, which compile
spatially resolved measurements of cold gas, ionized gas, stellar mass,
and SFR for a homogeneous sample of $\sim80$ nearby star-forming disk
galaxies \citep{2021ApJS..255...19L,2022AJ....164...43S,2023AJ....166..145L}.
The megatables tile each galaxy in matched apertures sampled on a regular
grid in galactocentric coordinates, and provide for each aperture aperture
photometry of CO(2--1) from \PHANGS-ALMA, \HI\ 21~cm from the \PHANGS-\HI\
companion programs, stellar mass surface density from environmental masks
of Spitzer/IRAC and WISE imaging, and multiple SFR tracers including
H$\alpha$, far-UV (FUV)+WISE-4 (W4) combinations, and an extinction-corrected
H$\alpha$+W4 estimator \citep{2022AJ....164...43S,2023ApJ...945L..19S}.

The megatables are released in two complementary aperture definitions:
hexagonal apertures, which provide a regular tessellation useful for
mapping context and for visual inspection, and two-dimensional
Gaussian-weighted apertures, which provide smoothly tapered measurements
that mitigate edge effects in the aperture sum and are the canonical
aperture for the \PHANGS\ pressure and SFR scaling analyses
\citep{2020ApJ...892..148S,2023ApJ...945L..19S}. \CGK{We adopt the
Gaussian-weighted apertures as canonical throughout this work and use the
hexagonal apertures only for context plots --- please confirm.}

For each galaxy, the gas, stellar, and SFR megatables are released as
separate FITS tables indexed by galaxy and aperture. We perform an inner
join on (galaxy, aperture) across the three tables, retaining only
apertures with valid entries in all joined columns. This yields a single
matched table of approximately $4.6\times10^{4}$ apertures across
$\sim80$ \CGK{galaxies --- please verify the final number after quality
cuts}, which serves as the starting point for the design-field and
validation-field selections described below.

\subsection{Derived Quantities}
\label{sec:derived}

The five environmental parameters that enter \TIGRESS\ as initial
conditions form the design-field vector
\begin{equation}
  \thetaenv \;=\; \left(\Sgas,\, \Sstar,\, \Hstar,\, \Omrot,\, \qshear\right).
  \label{eq:thetaenv}
\end{equation}
We construct each component as follows.

\paragraph{Gas surface density.}
The total neutral gas surface density is the sum of atomic and molecular
contributions,
\begin{equation}
  \Sgas \;=\; \Satom + \Smol,
\end{equation}
where $\Satom$ is derived from the \HI\ 21~cm aperture photometry and
$\Smol$ is derived from the \PHANGS-ALMA CO(2--1) aperture photometry
assuming the metallicity-dependent CO-to-H$_2$ conversion factor adopted
by \citet{2022AJ....164...43S}. Both components include the standard
factor of 1.36 to account for helium. The molecular fraction
$\fmol = \Smol/\Sgas$ is computed from the same components and is used
only as a validation quantity (see below); it is not part of $\thetaenv$.

\paragraph{Stellar surface density.}
The stellar surface density $\Sstar$ is taken from the \PHANGS\ stellar
mass maps based on Spitzer/IRAC 3.6~$\mu$m or WISE-1 3.4~$\mu$m imaging,
calibrated to mass-to-light ratios using independent color information,
following the procedure described in \citet{2021ApJS..255...19L} and
\citet{2025ApJ...989...66V}.

\paragraph{Orbital frequency.}
The orbital angular frequency at each aperture is computed from the
circular velocity of the host galaxy's rotation curve,
\begin{equation}
  \Omrot \;=\; \frac{V_{\rm circ}(R)}{R},
\end{equation}
where $R$ is the galactocentric radius of the aperture center and
$V_{\rm circ}(R)$ is the smooth axisymmetric rotation curve adopted by the
\PHANGS\ dynamical modeling pipeline \citep{2025ApJ...989...66V}.

\paragraph{Stellar scale height.}
The stellar scale height $\Hstar$ is not directly observed and must be
inferred from a stellar dynamical model. We adopt the values tabulated by
\citet{2025ApJ...989...66V}, who fit an axisymmetric stellar mass model
constrained by the rotation curve and the observed stellar surface density
profile, and report the implied vertical scale height under the assumption
that the stellar disk is locally flattened with an axis ratio
characteristic of nearby disk galaxies. \CGK{Please confirm whether
the $\Hstar$ used here is the exponential scale height or the
$\mathrm{sech}^{2}$ scale height in the \citet{2025ApJ...989...66V}
convention; this matters for the mapping to $\zstar$ in
\autoref{sec:tigress_mapping}.}

\paragraph{Shear parameter.}
The shear parameter quantifies the local logarithmic gradient of the
rotation curve and is defined as
\begin{equation}
  \qshear \;\equiv\; -\frac{d\ln \Omrot}{d\ln R}
        \;=\; 1 - \frac{d\ln V_{\rm circ}}{d\ln R}.
  \label{eq:qshear}
\end{equation}
For a flat rotation curve $\qshear = 1$, for solid-body rotation
$\qshear = 0$, and for Keplerian rotation $\qshear = 1.5$. We compute
$\qshear$ aperture-by-aperture from the numerical derivative of the
\PHANGS\ rotation curves at the galactocentric radius of each aperture.
Disk-dominated galaxies are physically bounded by $0 < \qshear \leq 1.5$;
this bound is imposed at the quality-cut stage (\autoref{sec:quality_cuts}).

\subsection{Quality Cuts and Completeness}
\label{sec:quality_cuts}

The reference set $\Dref$ used for sampling is defined by validity cuts
on the five design fields only:
\begin{equation}
  \Dref \;=\; \bigl\{\, \thetaenv \;\big|\;
              \Sgas,\, \Sstar,\, \Hstar,\, \Omrot > 0;\;
              0 < \qshear \leq 1.5 \,\bigr\}.
  \label{eq:dref}
\end{equation}
The positivity requirements on $\Sgas$, $\Sstar$, $\Hstar$, and $\Omrot$
remove apertures with missing or non-finite values in any of the four
quantities. The upper bound $\qshear \leq 1.5$ excludes apertures whose
locally fit rotation-curve derivative is super-Keplerian; such cases occur
near rotation-curve transitions or in noisy outer-disk apertures and are
unphysical for the smooth axisymmetric models that anchor the
\TIGRESS\ initial conditions.

Crucially, we do \emph{not} impose detection cuts on $\fmol$, on the
\HI--CO partition $(\Satom,\Smol)$, or on any of the SFR tracers when
defining $\Dref$. These quantities have their own, generally narrower,
completeness masks: CO is undetected at the lowest gas surface densities,
H$\alpha$ is suppressed in regions of low recent star formation, and W4
saturates in the brightest disks. If validation-field completeness were
imposed on the design selection, the \PHANGS\ prior would be biased toward
high-detection environments---preferentially CO-bright, H$\alpha$-bright,
inner-disk apertures---and the resulting simulation suite would
under-sample the low-pressure, atomic-dominated outer disks that form a
substantial fraction of the total \PHANGS\ disk area. By restricting the
selection cuts to the design fields, the prior reflects the full
disk-area-weighted distribution of galactic environments, and the
validation quantities can be compared against \TIGRESS\ outputs without
introducing a selection bias \CGK{(this argument should be verified against
the actual \PHANGS\ completeness statistics --- please cross-check).}

Validation against \PHANGS\ is performed using only those apertures that
additionally pass the relevant validation-field completeness masks
(detected CO for $\fmol$ and $\Smol$, detected SFR tracer for $\Ssfr$,
and so forth). These masks are applied per-quantity at the validation
stage and are not propagated back into $\Dref$. The full set of design
and validation fields is summarized in \autoref{tab:fields}.

\begin{deluxetable*}{lllp{0.45\textwidth}}
\tablecaption{Design and validation fields used in this work.\label{tab:fields}}
\tablehead{
  \colhead{Field} & \colhead{Symbol} & \colhead{Unit} & \colhead{Role}
}
\startdata
\multicolumn{4}{c}{\textit{Design fields} $\thetaenv$ (selection of $\Dref$)} \\
\hline
Gas surface density        & $\Sgas$  & $\Surf$              & TIGRESS input \\
Stellar surface density    & $\Sstar$ & $\Surf$              & TIGRESS input \\
Stellar scale height       & $\Hstar$ & $\pc$                & TIGRESS input (mapped to $\zstar$) \\
Orbital frequency          & $\Omrot$ & $\kms\,\kpc^{-1}$    & TIGRESS input \\
Shear parameter            & $\qshear$ & dimensionless, $0<\qshear\leq 1.5$ & TIGRESS input \\
\hline
\multicolumn{4}{c}{\textit{Validation fields} (not used for selection of $\Dref$)} \\
\hline
Molecular fraction         & $\fmol$  & dimensionless, $0\leq \fmol \leq 1$ & gas-partition validation \\
Atomic surface density     & $\Satom$ & $\Surf$              & gas-partition validation \\
Molecular surface density  & $\Smol$  & $\Surf$              & gas-partition validation \\
SFR surface density        & $\Ssfr$  & $\sfrunit$           & SFR validation \\
Dynamical equilibrium pressure & $\PDE$ & $\Punit$           & \PRFM\ validation \\
Effective velocity dispersion & $\seff$ & $\kms$             & \PRFM\ validation \\
\enddata
\tablecomments{Design fields are subject only to the validity cuts in
\autoref{eq:dref} and define the reference set $\Dref$ used by the
sampling design. Validation fields carry separate per-quantity
completeness masks and are compared against emergent \TIGRESS\ outputs
in \autoref{sec:validation}.}
\end{deluxetable*}

\section{Sampling Design}
\label{sec:sampling}

This section describes the construction of an observationally informed prior
over the environmental parameters that enter \TIGRESS\ as initial conditions,
and the quasi-random design that converts this prior into a concrete set of
simulation inputs. The design fields, their identification with \PHANGS\
aperture measurements, and the validation fields reserved for post-simulation
comparison are defined in \autoref{sec:data}; the mapping to \TIGRESS\ initial
conditions is detailed in \autoref{sec:tigress_mapping}.

%-----------------------------------------------------------------
\subsection{Statistical Goal}
\label{sec:sampling:goal}
%-----------------------------------------------------------------

The objective of the sampling procedure is to construct an informed prior over
the environmental design vector
\begin{equation}
\thetaenv \;=\; \rbrackets{\Sgas,\ \Sstar,\ \Hstar,\ \Omrot,\ \qshear},
\label{eq:theta_env}
\end{equation}
conditioned on a target gas surface density band,
\begin{equation}
p\rbrackets{\thetaenv \,\big|\, \log\Sgas \in
\sbrackets{\log\Sgas^{\rm tgt} \pm \Delta}}.
\label{eq:prior_definition}
\end{equation}
This prior is \emph{not} an exact resampling of the \PHANGS\ aperture
distribution. Exact resampling would inherit the survey's selection effects
(e.g.\ inclination, distance, sensitivity, aperture footprint) and would tile
parameter space inefficiently for emulator and simulation-based inference
(\SBI) training. Instead, the goal is a compact set of simulation inputs that:
(i) preserves the dominant observed correlations among the five design fields;
(ii) covers the relevant parameter volume sufficiently densely for emulator
training and posterior coverage tests; and (iii) excludes parameter
combinations that are unphysical or implausible.

We deliberately exclude the molecular gas fraction $\fmol$ and the star
formation surface density $\Ssfr$ from the design vector. Neither is a
\TIGRESS\ input parameter: in the \TIGRESS\ framework both are emergent
outputs of the explicit thermochemistry and self-consistent star formation
prescription \citep{2024ApJ...972...67K}.\CGK{Add a second TIGRESS-NCR reference here if desired (e.g.\ Kim et al.\ 2023 ApJS thermochemistry paper); the key 2023ApJS..264...10K is not yet in references.bib and should be verified in ADS before insertion.} Requiring valid
$\fmol$ and $\Ssfr$ measurements in the reference set would furthermore bias
the prior toward high-completeness environments in which CO and SFR tracers
are simultaneously detected, suppressing the very low-$\fmol$ and
low-$\Ssfr$ regions that the simulation suite must cover. These two fields
are instead retained as \emph{validation quantities} to be compared against
the corresponding emergent distributions from the completed simulation
suite (\autoref{sec:results}).

%-----------------------------------------------------------------
\subsection{Design Fields and Validation Fields}
\label{sec:sampling:fields}
%-----------------------------------------------------------------

We partition the relevant \PHANGS\ aperture-level quantities into two
disjoint sets: design fields, which serve as \TIGRESS\ initial conditions
and over which the prior is constructed, and validation fields, which are
emergent in \TIGRESS\ and used only for post-simulation comparison.

\paragraph{Design fields.}
The five-dimensional design vector $\thetaenv$ (defined in
\autoref{eq:theta_env}) contains the gas surface density $\Sgas$, the
stellar surface density $\Sstar$, the stellar scale height $\Hstar$,
the orbital frequency $\Omrot$, and the shear parameter
$\qshear \equiv -d\ln\Omrot/d\ln R$.
These five quantities are sufficient to specify the local shearing-box
\TIGRESS\ initial conditions up to the dark matter midplane density
$\rhodm$, which is derived from the rotation curve as described in
\autoref{sec:tigress_mapping}.

\paragraph{Validation fields.}
The validation set comprises emergent quantities that are diagnostic of
the multiphase ISM and star formation, but are not used to specify the
simulations:
\begin{equation}
\thetaenv^{\rm val} = \rbrackets{\fmol,\ \Satom,\ \Smol,\ \Ssfr,\ \PDE,\ \seff}.
\end{equation}
Among these, $\PDE$ and $\seff$ are derived quantities in the \PRFM\
framework \citep{2022ApJ...936..137O} that are computed from the design
fields via the vertical equilibrium and feedback closure relations;
their \PHANGS\ values therefore depend on the same five design fields
plus the choice of \PRFM\ effective velocity dispersion model. We use
the design-field/validation-field split exclusively to define which
distributions are matched in the prior construction (design fields)
versus checked in post-simulation comparison (validation fields).

%-----------------------------------------------------------------
\subsection{Reference Pixel Selection}
\label{sec:sampling:reference}
%-----------------------------------------------------------------

The reference dataset $\Dref$ from which the prior is fit is the subset of
\PHANGS\ apertures \citep{2022AJ....164...43S,2023ApJ...945L..19S} that
satisfy a target $\Sgas$ band and basic physical validity. Specifically,
a \PHANGS\ aperture $i$ is included in $\Dref$ if
\begin{align}
\bigl|\log\Sgas^{(i)} - \log\Sgas^{\rm tgt}\bigr| &\leq \Delta,
\label{eq:sgas_band}\\
\Sgas^{(i)},\ \Sstar^{(i)},\ \Hstar^{(i)},\ \Omrot^{(i)} &> 0,
\label{eq:positivity}\\
0 &< \qshear^{(i)} \leq 1.5.
\label{eq:q_bound}
\end{align}
The band width $\Delta$ defines the local extent of $\Dref$ in $\log\Sgas$
and is chosen to balance two competing requirements: $\Delta$ must be wide
enough to provide a statistically meaningful sample for the kernel density
estimator, yet narrow enough that the conditional joint distribution of the
remaining design fields is well approximated as locally stationary in
$\log\Sgas$. A fiducial choice $\Delta = 0.2\dex$ is adopted in this work
and discussed further in \autoref{sec:results}.

The upper bound on the shear parameter, $\qshear \leq 1.5$, follows from
the definition $\qshear = -d\ln\Omrot/d\ln R$: $\qshear=0$ corresponds to
solid-body rotation, $\qshear=1$ to a flat rotation curve, and
$\qshear=1.5$ to Keplerian rotation around a point mass. Values
$\qshear > 1.5$ would correspond to super-Keplerian rotation curves that
cannot be produced by a positive enclosed mass distribution and are
therefore excluded as unphysical. We similarly exclude $\qshear \leq 0$,
which corresponds to rotation curves that increase faster than linearly
with radius and is not a regime in which the local shearing-box
approximation underlying \TIGRESS\ is expected to apply.

Crucially, validity of $\fmol$, $\Ssfr$, or any other tracer-specific
quantity is \emph{not} required for inclusion in $\Dref$. \CGK{This choice
maximizes the number of usable reference pixels — particularly in the
low-$\Sgas$, low-$\fmol$ outer-disk regime where CO is undetected and the
PHANGS megatable returns upper limits on $\Smol$ — but should be verified
against the megatable flag convention to ensure that apertures with all
five design fields validly measured are retained.}

%-----------------------------------------------------------------
\subsection{Kernel Density Estimator Prior}
\label{sec:sampling:kde}
%-----------------------------------------------------------------

We work throughout in log-transformed design space. Define the
five-dimensional log-vector
\begin{equation}
\bm{w} \;\equiv\;
\rbrackets{\log\Sgas,\ \log\Sstar,\ \log\Hstar,\ \log\Omrot,\ \log\qshear},
\label{eq:w_def}
\end{equation}
so that the reference set in log-space is $\cbrackets{\bm{w}_i}_{i=1}^{N}$
with $N = |\Dref|$. The log transformation is natural because each design
field spans more than an order of magnitude across \PHANGS\ and is
approximately log-normally distributed within fixed $\Sgas$ bands; it also
makes the KDE bandwidth selection scale-invariant under multiplicative
rescaling of any field.

We fit a Gaussian kernel density estimator to the reference log-vectors,
\begin{equation}
\hat{p}(\bm{w}) \;=\; \frac{1}{N}\,\sum_{i=1}^{N}
K_{\bm{H}}\rbrackets{\bm{w} - \bm{w}_i},
\label{eq:kde}
\end{equation}
where $K_{\bm{H}}$ is the multivariate Gaussian kernel
\begin{equation}
K_{\bm{H}}(\bm{u}) \;=\;
\rbrackets{2\pi}^{-5/2}\,\det(\bm{H})^{-1/2}\,
\exp\!\sbrackets{-\tfrac{1}{2}\,\bm{u}^{\top}\bm{H}^{-1}\bm{u}},
\label{eq:gaussian_kernel}
\end{equation}
and $\bm{H}$ is a symmetric positive-definite bandwidth matrix. We adopt
the multivariate Scott's rule for the full bandwidth matrix,
\begin{equation}
\bm{H} \;=\; \left(\frac{4}{d+2}\right)^{\!\!2/(d+4)}
\!\! N^{-2/(d+4)}\,\hat{\bm{\Sigma}},
\label{eq:scott_rule}
\end{equation}
where $d = 5$ is the number of design dimensions,
$N = |\Dref|$ is the reference sample size, and $\hat{\bm{\Sigma}}$ is
the empirical covariance of the $N$ reference log-vectors; for $d=5$
this gives $\bm{H} = (4/7)^{2/9} N^{-2/9} \hat{\bm{\Sigma}}$.
Multiplying the full empirical covariance by a scalar factor ensures
that $\bm{H}$ is positive definite and correctly encodes the observed
correlations among all five log-fields. In practice we will compare
against likelihood-based cross-validation to verify that this bandwidth
does not oversmooth the joint distribution.
Because all five design fields are strictly positive in $\Dref$ and
their upper bounds (most importantly $\qshear \leq 1.5$) are handled
by the post-mapping rejection step rather than by a boundary-corrected
KDE, the KDE in log-space requires no special boundary treatment.
This is in contrast to designs that include the molecular fraction
$\fmol \in [0,1]$, where a logit transform is needed to avoid
density leakage outside the physical support.

From $\hat{p}(\bm{w})$ we draw a large auxiliary KDE sample of size
$M \gg N$ (we use $M = 10^{5}$) and from this sample compute, for each
design coordinate $j = 1,\ldots,5$, the marginal cumulative distribution
function $F_j(w_j)$ and its inverse, the quantile function
\begin{equation}
Q_j(u) \;=\; F_j^{-1}(u), \qquad u \in (0,1).
\label{eq:quantile_function}
\end{equation}
The quantile functions are tabulated on a dense grid in $u$ and used as
look-up tables in the Sobol mapping (\autoref{sec:sampling:sobol}).
We emphasize that $\hat{p}(\bm{w})$ retains the full multivariate
correlation structure through the off-diagonal entries of $\bm{H}$,
but the marginal quantile functions $\cbrackets{Q_j}$ encode only the
one-dimensional marginals of $\hat{p}$. The dimension-wise
inverse-CDF mapping in \autoref{eq:sobol_map} therefore produces a
sample whose one-dimensional marginal distributions match those of
$\hat{p}$ \emph{exactly}, while the joint distribution corresponds to
an independence copula applied to these marginals. The joint
correlation structure of $\hat{p}$ is thus reproduced only
approximately, with fidelity depending on how well an independence
copula approximates the true KDE joint. In practice this approximation
is adequate when the pairwise correlations among the log design fields
are moderate, as is typical within a fixed $\Sgas$ band; we assess the
quality of the joint match quantitatively in \autoref{sec:results}.

%-----------------------------------------------------------------
\subsection{Sobol Sequence Design}
\label{sec:sampling:sobol}
%-----------------------------------------------------------------

Having reduced the prior to its marginal quantile functions on a unit
hypercube, we map a quasi-random low-discrepancy sequence through these
functions to produce the simulation design. We adopt a scrambled Sobol
sequence \citep{Sobol1967,Owen1998} rather than a pseudo-random or Latin
Hypercube Sampling (LHS) design for two reasons. First, for any finite
sample size $n$, Sobol sequences achieve a star-discrepancy of
$O(n^{-1}(\log n)^{d})$ in $d$ dimensions, asymptotically smaller than
the $O(n^{-1/2})$ scaling of pseudo-random samples and uniformly better
than LHS in projections onto coordinate subspaces. Second, and more
importantly for an incrementally grown simulation suite, the Sobol
construction has a nested progressive property: the first $2^{k}$ points
of the sequence form a balanced design at all $k$, so doubling the sample
size from $2^{k}$ to $2^{k+1}$ adds new design points without displacing
any of the original ones. LHS, by contrast, must in general be
re-stratified — and therefore re-drawn from scratch — whenever the
sample size changes.

\paragraph{Sampling algorithm.}
Given a target sample size $n$, the design is generated as follows.
\begin{enumerate}
\item Draw a scrambled Sobol sequence
$\cbrackets{\bm{u}_k}_{k=1}^{n} \subset [0,1]^{5}$ using
\texttt{scipy.stats.qmc.Sobol(d=5, scramble=True)} with a fixed random
seed for reproducibility.
\item Map each Sobol point through the marginal quantile functions of
\autoref{eq:quantile_function} to obtain log-space design points,
\begin{equation}
w_{k,j} \;=\; Q_j\rbrackets{u_{k,j}},
\qquad j = 1,\ldots,5.
\label{eq:sobol_map}
\end{equation}
\item Convert to physical-space design points,
\begin{equation}
\theta_{k,j} \;=\; 10^{\,w_{k,j}},
\label{eq:to_physical}
\end{equation}
yielding $\bm{\theta}_k = (\Sgas, \Sstar, \Hstar, \Omrot, \qshear)_k$.
\item For each design point with $\qshear_k > 1.5$, draw the next
un-used Sobol point from the sequence and repeat steps 2--3 until an
accepted point is obtained. Let $n_{\rm extra}$ denote the total
number of rejected-and-replaced draws. The accepted set of $n$ points
is formed from the first $n + n_{\rm extra}$ Sobol indices, with the
$n_{\rm extra}$ rejected indices recorded so that future extensions
begin from index $n + n_{\rm extra} + 1$. This bookkeeping preserves
the progressive nesting property (\autoref{sec:sampling:progressive}).
\end{enumerate}
The shear bound is the only post-mapping rejection criterion; positivity
is automatically guaranteed by the log-space construction. Because the
$\qshear$ marginal in $\Dref$ is already upper-truncated by
\autoref{eq:q_bound}, the rejection fraction $n_{\rm extra}/n$ is
expected to be small; we report the realized value in
\autoref{sec:results}.

Sample sizes that are powers of two, $n = 2^{m}$ with $m$ a non-negative
integer, are preferred because the Sobol sequence's discrepancy bounds
are sharp at these sizes and the construction is balanced in projections
to dyadic subintervals. Arbitrary $n$ is permitted at the cost of a mild
loss of projection uniformity.

%-----------------------------------------------------------------
\subsection{Progressive Extension of the Suite}
\label{sec:sampling:progressive}
%-----------------------------------------------------------------

A central operational advantage of the Sobol-based design is that the
suite is extensible at any later stage without re-drawing the existing
simulations. Let $S(n)$ denote the set of accepted design points produced by the
algorithm of \autoref{sec:sampling:sobol} at sample size $n$.
The underlying Sobol sequence has the nesting property: the first
$2^{k}$ sequence points are a subset of the first $2^{k+1}$ sequence
points. Provided the same random seed and the same physical rejection
criterion ($\qshear > 1.5$) are applied at every stage, the
accepted-point sets inherit this property,
\begin{equation}
S(2^{k}) \;\subset\; S(2^{k+1}) \;\subset\; S(2^{k+2}) \;\subset\; \cdots,
\label{eq:sobol_nesting}
\end{equation}
so that any simulation completed at stage $k$ is automatically a member
of the design at every subsequent stage. In practice we generate
$n + n_{\rm extra}$ Sobol points at each stage, accept the first $n$
points that pass the rejection criterion, and record the consumed Sobol
indices so that future extensions begin from the next un-used point in
the same sequence. The first generation of the
suite targets $n \in \{64, 128, 256\}$ ($m \in \{6, 7, 8\}$),
corresponding to a sequence of doublings; this progression is set by the
available computational allocation rather than any statistical
threshold, and the discrepancy of the design improves monotonically
along the sequence.

At each stage, the new $2^{k}$ points fill the largest remaining gaps in
the previous-stage coverage of the unit hypercube and, by the inverse-CDF
mapping of \autoref{eq:sobol_map}, the largest remaining gaps in
log-design space weighted by the KDE prior. This contrasts sharply with
LHS-based designs, which must be re-stratified to add new points and
therefore either displace the existing design or accept a strictly
sub-optimal stratification on the enlarged design.

This progressive structure is closely analogous to the staged extension
strategy adopted by the \CAMELS\ project
\citep{2021ApJ...915...71V}, in which cosmological hydrodynamic
simulation suites have been incrementally enlarged as new scientific
questions are identified. The Sobol nesting provides for the
\PHANGS-informed \TIGRESS\ suite the same operational property that
\CAMELS\ achieves through its explicit batch design: completed
simulations are never invalidated by the addition of new ones.

%-----------------------------------------------------------------
\subsection{Block-Structured Extension to Physics Parameters}
\label{sec:sampling:physics}
%-----------------------------------------------------------------

The first-generation suite holds all physics parameters fixed at fiducial
solar-neighborhood values (\autoref{sec:tigress_mapping}). Subsequent
generations will vary uncertain physics parameters such as the gas-phase
metallicity $\Zg$, the dust-to-gas ratio $\Zd$, the cosmic-ray
ionization rate $\xicr$, the assumed initial mass function, and
parameters of the supernova feedback prescription. We collect these
into a $d_{\rm phys}$-dimensional physics design vector $\thetaphys$.

We adopt a block-structured Sobol design over $(\thetaenv, \thetaphys)$,
defined as follows. Let
\begin{equation}
\cbrackets{\bm{u}^{\rm env}_k}_{k=1}^{n_{\rm env}}
\;\subset\; [0,1]^{5},
\qquad
\cbrackets{\bm{u}^{\rm phys}_\ell}_{\ell=1}^{n_{\rm phys}}
\;\subset\; [0,1]^{d_{\rm phys}}
\label{eq:two_sobol}
\end{equation}
be two independent scrambled Sobol sequences, one over the environmental
block and one over the physics block. The combined design is constructed by index-aligned pairing:
simulation $k$ uses environmental point $\bm{u}^{\rm env}_k$ and
physics point $\bm{u}^{\rm phys}_k$, so the total number of
simulations equals $\min(n_{\rm env}, n_{\rm phys})$. Each block is
mapped through its respective prior: $\bm{u}^{\rm env}_k$ through the
KDE-based quantile functions of \autoref{eq:sobol_map}, and
$\bm{u}^{\rm phys}_k$ through prior quantile functions appropriate to
each physics parameter (e.g.\ log-uniform in $\Zg$ over a chosen range).
When the two blocks differ in size, the smaller block is extended by
drawing additional Sobol points before pairing.

The block construction has two advantages. First, the two blocks may be
extended independently: doubling $n_{\rm env}$ at fixed $n_{\rm phys}$
yields a denser environmental sampling at the original set of physics
points, while doubling $n_{\rm phys}$ at fixed $n_{\rm env}$ adds physics
variations to the existing environmental design. Both extensions inherit
the nesting property of \autoref{eq:sobol_nesting} within their
respective blocks. Second, the two priors are kept manifestly
independent, which is desirable when the physics-parameter prior is
itself uncertain or under active revision: an update to $p(\thetaphys)$
does not require re-drawing the environmental design.

This block structure is analogous to the separation between
``cosmological'' and ``astrophysical'' parameter suites in the \CAMELS\
program \citep{2021ApJ...915...71V}, which treats the two parameter
groups as independent design axes that can be extended along either axis
without re-running the other.

\paragraph{Option A: pre-allocated joint design.}
If the full set of varied physics parameters and their priors are known
\emph{a priori}, an alternative is to draw a single
$(5+d_{\rm phys})$-dimensional Sobol sequence over the joint space and
map each block through its respective prior. This yields the optimal
joint discrepancy at all sample sizes and preserves the nesting property
of \autoref{eq:sobol_nesting} in the joint space, but at the cost of
fixing $d_{\rm phys}$ from the outset: adding a new physics parameter
later requires re-drawing the design. We therefore prefer the
block-structured construction for the staged suite considered here, and
reserve the pre-allocated joint design for a future generation in which
the physics-parameter list is finalized.

%-----------------------------------------------------------------
\subsection{Expanded Prior for Rare Environments}
\label{sec:sampling:expanded}
%-----------------------------------------------------------------

The KDE prior of \autoref{sec:sampling:kde} faithfully reproduces the
\PHANGS\ design-field distribution within the target $\Sgas$ band, but
by construction it under-samples regions of parameter space that are
under-represented in \PHANGS. Such regions include, for example, very
low orbital frequencies $\Omrot$ characteristic of low-surface-brightness
or outer-disk environments, large stellar scale heights $\Hstar$
characteristic of thickened or perturbed disks, and shear parameters
$\qshear$ near the rising-rotation-curve ($\qshear \to 0$) or Keplerian
($\qshear \to 1.5$) extremes. For emulator and \SBI\ training it is
desirable to extend coverage into these tails to guard against
extrapolation in posterior inference.

We implement such extension in two complementary ways. (i) Broadening
the KDE bandwidth, $\bm{H} \to s\,\bm{H}$ with $s > 1$, inflates the
prior in all directions while preserving the empirical mean and
correlation structure; this is sufficient for modest extrapolation.
(ii) Directional extension of the acceptance bounds in selected
log-space directions, implemented by augmenting the quantile functions
$Q_j(u)$ at the tails (i.e.\ extending $Q_j$ for $u$ near $0$ or $1$
beyond the support of $\Dref$), supports targeted extension into rare
regimes such as low-$\Omrot$ environments. Expanded-prior simulations
are flagged at draw time and analyzed separately from the
\PHANGS-prior simulations in the fairness diagnostics
(\autoref{sec:sampling:fairness}) and downstream emulator validation.
\CGK{The precise broadening factor $s$ and the directional extensions
applied in the first-generation suite will be reported in
\autoref{sec:results}; the choice of these is closely tied to which
PHANGS-poor environments the inference program ultimately targets.}

%-----------------------------------------------------------------
\subsection{Sampling Fairness Diagnostic}
\label{sec:sampling:fairness}
%-----------------------------------------------------------------

We define a quantitative diagnostic of how faithfully the simulation
design reproduces the \PHANGS\ design-field distribution within the
target $\Sgas$ band. For each design field $a \in \thetaenv$, let
$\hat{q}^{\rm sample}_{a}(\tau)$ and $\hat{q}^{\rm ref}_{a}(\tau)$
denote the empirical $\tau$-th quantile of $\log_{10} a$ in the
simulation design and in $\Dref$, respectively (we use $\tau$ for
quantile probability levels to distinguish from the KDE marginal
quantile functions $Q_j$ of \autoref{eq:quantile_function}).
The per-field quantile mismatch is
\begin{equation}
\varepsilon_a \;=\;
\max_{\tau \in \cbrackets{0.1,\,0.2,\,\ldots,\,0.9}}
\Bigl|\,\hat{q}^{\rm sample}_{a}(\tau)
- \hat{q}^{\rm ref}_{a}(\tau)\,\Bigr|,
\label{eq:eps_a}
\end{equation}
in units of dex, and the overall design fairness is the worst-field
mismatch,
\begin{equation}
\varepsilon_{\rm env} \;=\;
\max_{a \in \thetaenv}\,\varepsilon_a.
\label{eq:eps_env}
\end{equation}
A design is considered \emph{fair} with respect to $\Dref$ if
$\varepsilon_{\rm env} \leq \varepsilon_{\rm tol}$, where we adopt
$\varepsilon_{\rm tol} = 0.10$--$0.20\dex$ as a fiducial tolerance.
This range is motivated by the typical aperture-level measurement
uncertainty in the underlying \PHANGS\ design fields, which is
itself of order $\sim 0.1\dex$ for the surface densities and somewhat
larger for $\Hstar$; a design that matches the reference marginal
quantiles to better than the measurement uncertainty is not
meaningfully distinguishable from the reference within $\Dref$.
\CGK{The 0.10--0.20\,dex tolerance is a working choice; a
quantitative justification from the PHANGS measurement-uncertainty
budget should be added when those numbers are tabulated.}

The diagnostic in \autoref{eq:eps_env} applies exclusively to the
five \emph{design} fields. The validation fields $\fmol$, $\Satom$,
$\Smol$, $\Ssfr$, $\PDE$, and $\seff$ are not used to define the
prior, and they cannot be matched at design time because they are
not \TIGRESS\ inputs. Instead, the simulation outputs for the
validation fields will be compared against the corresponding
\PHANGS\ distributions in \autoref{sec:results} using analogous
quantile-mismatch metrics, providing a post-hoc verification that
the simulated ISM populates the same emergent regime as the
observed one. A successful program is one in which
$\varepsilon_{\rm env} \leq \varepsilon_{\rm tol}$ by design and the
corresponding validation-field mismatches are also small without
having been tuned.

\section{TIGRESS-NCR Parameter Mapping}
\label{sec:tigress_mapping}

\subsection{TIGRESS-NCR Input Parameters}

The \TIGRESS\ framework \citep{2024ApJ...972...67K} solves the equations of
ideal MHD with self-gravity and explicit non-equilibrium thermo-chemistry in a
local shearing-box geometry. The initial conditions are specified by six
parameters:
\begin{enumerate}
  \item $\Sgas$: total gas surface density [$\Surf$]
  \item $\Sstar$: stellar surface density [$\Surf$]
  \item $\zstar$: stellar scale height [pc]
  \item $\Omrot$: orbital frequency [km\,s$^{-1}$\,kpc$^{-1}$]
  \item $\qshear$: shear parameter $q \equiv -d\ln\Omrot/d\ln R$ [dimensionless]
  \item $\rhodm$: dark-matter volume density at the midplane [$\rhounit$]
\end{enumerate}

\subsection{Mapping from PHANGS Design Fields to TIGRESS-NCR}

Four of the five \PHANGS\ design fields map directly to \TIGRESS\ inputs with
no conversion required:
\begin{align}
(\Sgas,\, \Sstar,\, \Hstar,\, \Omrot,\, \qshear)_{\rm PHANGS}
&\nonumber\\
\longrightarrow\;
(\Sgas,\, \Sstar,\, \zstar,&\, \Omrot,\, \rhodm,\, \qshear)_{\rm TIGRESS}
\end{align}

\paragraph{Stellar scale height.}
The \PHANGS\ stellar scale height $\Hstar$ is derived as
\begin{equation}
  \Hstar \;=\; \frac{\Sstar}{4\,\rho_{\star,\rm mid}},
  \label{eq:hstar}
\end{equation}
where $\rho_{\star,\rm mid}$ is the stellar midplane volume density from the
axisymmetric mass model of \citet{2025ApJ...989...66V}. This is directly
equivalent to the \TIGRESS\ stellar scale height parameter:
\begin{equation}
  \zstar \;=\; \Hstar.
\end{equation}
No additional conversion is required.

\paragraph{Dark matter density.}
The dark-matter midplane density $\rhodm$ is not directly available from the
\PHANGS\ megatables and cannot be reliably derived aperture-by-aperture at the
$\sim$kpc resolution of this survey. For the first simulation suite we therefore
hold $\rhodm$ fixed at the fiducial value used in the \TIGRESS\ R8 model, which
corresponds to the solar-neighborhood environment at $\Sgas \approx 10\Surf$:
\begin{equation}
  \rhodm \;=\; 0.0064\rhounit.
  \label{eq:rhodm_fiducial}
\end{equation}
This value is consistent with the local dark-matter density inferred from the
Milky Way rotation curve in the vicinity of the solar circle
\citep{2024ApJ...972...67K}. Treating $\rhodm$ as a fixed parameter is a
reasonable first approximation because $\rhodm$ enters the vertical equilibrium
primarily through its contribution to the total gravitational weight at large
$|z|$; for the gas surface densities and disk environments sampled here,
variations in $\rhodm$ of order a factor of two have a subdominant effect on
$\Ssfr$ and $\PDE$ compared to variations in $\Sgas$ and $\Sstar$.

\subsection{Summary of Parameter Mapping}

\begin{deluxetable*}{llll}
\tablecaption{PHANGS design field to \TIGRESS-NCR\ input mapping
\label{tab:mapping}}
\tablehead{
  \colhead{PHANGS field} & \colhead{Unit} &
  \colhead{TIGRESS parameter} & \colhead{Notes}
}
\startdata
$\Sgas$   & $\Surf$                   & $\Sgas$   & direct \\
$\Sstar$  & $\Surf$                   & $\Sstar$  & direct \\
$\Hstar$  & pc                        & $\zstar$  & direct; $\Hstar = \Sstar/(4\,\rho_{\star,\rm mid})$ \\
$\Omrot$  & km\,s$^{-1}$\,kpc$^{-1}$ & $\Omrot$  & direct \\
$\qshear$ & ---                       & $\qshear$ & direct; $0 < q \leq 1.5$ \\
---       & ---                       & $\rhodm$  & fixed at $0.0064\rhounit$ (R8 fiducial)
\enddata
\tablecomments{The stellar scale height $\Hstar$ computed from \PHANGS\ is
  identical in definition to the \TIGRESS\ parameter $\zstar$ and requires no
  conversion. The dark-matter density $\rhodm$ is not available from \PHANGS\
  at aperture level; it is held fixed at the solar-neighborhood R8 value for
  the first simulation suite (\autoref{eq:rhodm_fiducial}).}
\end{deluxetable*}

\subsection{Fiducial Physics Parameters}

The first simulation suite holds all physics parameters at fiducial values:
metallicity $\Zg = 1$ (solar), dust-to-gas ratio $\Zd = 1$, and standard
TIGRESS-NCR feedback and population-synthesis prescriptions
\citep{2024ApJ...972...67K}. Extension to physics parameter variations
(metallicity, IMF, feedback yield) is described in \autoref{sec:discussion}.

\section{Results}
\label{sec:results}

\subsection{Sampling Diagnostics}

Figure~\ref{fig:field_distributions} shows the marginal distributions of the five
design fields for the KDE-Sobol samples at $n = 64$ and $n = 128$, compared to the
\PHANGS\ reference set. The fairness metric $\epsilon_{\rm env}$
(\autoref{sec:fairness}) is reported for each sample size.

\begin{figure}[ht]
  \centering
  % \includegraphics[width=\columnwidth]{field_distributions.png}
  % Figure placeholder: replace with actual output from figures/phangs/
  \caption{Marginal distributions of the five environmental design fields
           ($\Sgas$, $\Sstar$, $\Hstar$, $\Omrot$, $\qshear$) for the
           KDE-Sobol samples ($n=64$: blue; $n=128$: orange) compared to
           the \PHANGS\ reference set (gray). Quantile fairness $\epsilon$
           is indicated for each field.
           \textit{[PLACEHOLDER --- figure to be updated with Sobol-based results]}}
  \label{fig:field_distributions}
\end{figure}

\subsection{Simulated ISM Properties}
\label{sec:sim_results}

\textbf{[PLACEHOLDER]} This section will present the TIGRESS-NCR simulation results
for the first suite of $n$ runs with PHANGS-informed environmental parameters at
fiducial metallicity $\Zg = 1$.

Expected outputs to be reported:
\begin{itemize}
  \item Time-averaged $\Ssfr$, $\fmol$, $\PDE$, $\seff$ for each simulation
  \item ISM phase structure and midplane pressure as functions of $\Sgas$, $\Sstar$, $\Omrot$
  \item PRFM feedback yield calibration from the suite
\end{itemize}

\subsection{Validation Against PHANGS}
\label{sec:validation}

\textbf{[PLACEHOLDER]} This section will compare the emergent simulation properties
against the observed \PHANGS\ distributions.

\begin{figure}[ht]
  \centering
  % \includegraphics[width=\columnwidth]{validation_sfr_fmol.pdf}
  \caption{Comparison of emergent $\Ssfr$ and $\fmol$ from the \TIGRESS\ suite
           against the observed \PHANGS\ distributions at matched $\Sgas$.
           \textit{[PLACEHOLDER --- results pending]}}
  \label{fig:validation}
\end{figure}

\section{Discussion and Future Directions}
\label{sec:discussion}

The simulation design presented in this paper is the first stage of a
broader program: an observationally anchored, progressively extensible
ensemble of resolved multiphase ISM simulations that supports systematic
comparison with \PHANGS\ and statistical inference of ISM physical
parameters. In this section we recapitulate the sampling design
(\autoref{sec:disc_summary}), outline how the block-structured Sobol
construction will be extended to physics parameters
(\autoref{sec:disc_phys}), describe the synthetic observation pipeline
that will close the loop with resolved \PHANGS\ maps
(\autoref{sec:disc_synth}), discuss summary statistics beyond one-point
aperture quantities (\autoref{sec:disc_twopoint}), and frame the
connection to the inference framework developed in the companion paper
(\autoref{sec:disc_paperII}).

\subsection{Summary of the Sampling Design}
\label{sec:disc_summary}

We have presented a \PHANGS-informed simulation design for the
\TIGRESS\ framework. Five environmental parameters---the gas surface
density $\Sgas$, stellar surface density $\Sstar$, stellar scale height
$\Hstar$, orbital frequency $\Omrot$, and shear parameter
$\qshear$---define the design vector $\thetaenv$ and enter \TIGRESS\
as initial and boundary conditions. The sampling proceeds in two steps:
a kernel density estimator (KDE) is fit to the joint \PHANGS\ aperture
distribution of $\thetaenv$ and used to define the marginal quantile
functions, and a scrambled Sobol quasi-random sequence is mapped through
these quantile functions to generate the simulation design.

Two features of this construction are essential for the program as a
whole. First, because the Sobol sequence is deterministic and
hierarchical, the suite can be extended progressively: doubling the
sample size adds new low-discrepancy points without displacing or
re-running any earlier simulation. Second, the molecular fraction
$\fmol$ and SFR surface density $\Ssfr$ are deliberately excluded from
the input design and reserved as validation quantities, to be compared
against emergent \TIGRESS\ outputs.
\CGK{This separation of design fields from validation fields is what
allows the suite to be used both as a calibration set (for \PRFM\
yields) and as a prior for simulation-based inference; flag for review
that this framing is consistent with the abstract.}

\subsection{Extension to Physics Parameter Variations}
\label{sec:disc_phys}

The first suite holds the physics parameters at fiducial
solar-neighborhood values: solar metallicity ($\Zg = 1$), solar
dust-to-gas ratio ($\Zd = 1$), a standard initial mass function (IMF),
and the fiducial \TIGRESS-NCR\ feedback prescription. This choice
isolates the environmental dependence of ISM properties from physics
variations, but it also means that the inferred posteriors from this
suite are conditioned on these fiducial physics assumptions. The
natural next step is to relax this conditioning by introducing an
independent block of simulations that varies a vector of physics
parameters $\thetaphys$.

We anticipate four parameters as the principal targets of this
extension. First, the gas-phase metallicity $\Zg \in [0.1, 3]$ in solar
units, motivated by the \TIGRESS-NCR\ metallicity suite of
\citet{2024ApJ...972...67K}, who showed that $\Ssfr$ scales weakly with
metallicity at fixed environment, approximately as $\Ssfr \propto
\Zg^{0.3}$. \CGK{Verify the exponent and the sign; I am recalling this
from the K24 abstract but please confirm against the paper before this
text is finalized.} Second, the dust-to-gas ratio $\Zd$, which is
strongly but not perfectly correlated with $\Zg$ in observed galaxies
and which controls photoelectric heating, $\HH$ formation, and
radiation transfer independently of gas-phase line cooling. Third, the
upper-mass cutoff of the IMF, which sets the ionizing photon budget and
the supernova rate per unit star formation. Fourth, an overall feedback
yield normalization that re-weights the thermal versus turbulent
partitioning of $\Ytot$ relative to its calibrated \TIGRESS-NCR\ value;
this last parameter is included not because we expect the calibration
to be wrong, but to allow the inference framework to detect
inconsistencies between the simulated and observed yield decompositions.

The block-structured Sobol construction makes this extension natural.
Block~1 is the present environmental suite, a Sobol design of dimension
$d_{\rm env} = 5$ over $\thetaenv$ at fiducial physics. Block~2 will be
a Sobol design of dimension $d_{\rm phys}$ over $\thetaphys$, run at a
representative subset of environmental anchor points selected from
Block~1. The joint suite supports inference over the concatenated
parameter vector $(\thetaenv, \thetaphys)$ by combining the marginal
information from each block under an assumed factorization of the
prior. This mirrors the multi-suite philosophy of the \CAMELS\ project
\citep{2021ApJ...915...71V}, in which one-parameter variation suites,
Latin-hypercube suites, and extension suites jointly cover the
cosmology and astrophysics parameter space at controlled cost.
\CGK{Verify that this is a faithful characterization of the CAMELS suite
structure; the high-level analogy is correct but the suite labels and
their exact roles should be checked.}

\subsection{Synthetic Observation Pipeline}
\label{sec:disc_synth}

A direct comparison between \TIGRESS\ and \PHANGS\ at the level of
integrated aperture quantities, as performed in
\autoref{sec:results}, is a necessary but not sufficient test of the
simulation suite. The same aperture-averaged $\Ssfr$, $\fmol$, $\PDE$,
and $\seff$ can be produced by simulations that differ substantially in
the spatial structure of their gas and young stellar populations. To
test the simulations at the level of resolved morphology, and to
construct simulated observables that can be ingested by the inference
framework of the companion paper, we plan to develop a forward-modeling
pipeline that converts \TIGRESS\ snapshots into mock multi-wavelength
maps matched to the \PHANGS\ resolution and noise properties.

The pipeline will produce five primary tracers. (i)~Mid-infrared dust
emission, modeled from the \TIGRESS\ dust column and ambient radiation
field to produce synthetic JWST/MIRI and \emph{Spitzer}/MIPS $24\,\mu$m
maps. (ii)~Ionized-gas H$\alpha$ emission, computed from the \HII\
region catalog of each \TIGRESS\ snapshot with extinction applied
self-consistently from the simulated dust column. (iii)~Far-ultraviolet
continuum, computed from the young stellar population with the same
extinction treatment. (iv)~\HI\ 21\,cm column, derived from the neutral
atomic phase of \TIGRESS-NCR. (v)~CO line emission, computed either
from a calibrated CO-to-H$_2$ conversion as a function of metallicity
and radiation field, or from direct radiative transfer post-processing
of the \HH\ phase. \CGK{Decide which CO route to adopt; both have a
literature trail and the choice will need to be justified explicitly.}
Each tracer will then be convolved with the appropriate beam, projected
into the observational aperture grid, and resampled to the \PHANGS\
sensitivity and noise level.

This synthetic observation pipeline serves three purposes. It allows
direct comparison of \TIGRESS\ outputs to observed \PHANGS\ maps in the
observed plane, removing the modeling assumptions implicit in
converting observed maps to physical surface densities. It produces the
maps required as inputs to field-level inference (\autoref{sec:disc_paperII}).
And it provides a self-consistent training set for emulators that
predict observable maps directly from $\thetaenv$, bypassing the
intermediate step of physical-parameter estimation.

\subsection{Beyond One-Point Statistics}
\label{sec:disc_twopoint}

The validation and inference targets used in this paper are one-point
aperture summaries: the mean $\Ssfr$, $\fmol$, $\PDE$, $\seff$, and
feedback yield $\Ytot$ within a fixed aperture. These quantities are
informative about the pressure balance and the mean efficiency of
feedback, but they discard the spatial coherence of star formation and
gas distribution that is plausibly the most discriminating signature
between feedback models with similar globally averaged behavior.

A natural first extension is to two-point statistics: the angular power
spectrum or two-point correlation function of $\Sgas$, $\Ssfr$, and
$\PDE$ maps within each aperture or across each simulated patch.
Two-point statistics constrain the spatial coherence length of star
formation, the slope of the turbulent cascade in the cold gas, and the
relative phasing of gas and SFR tracers. \CGK{Mention specific
turbulent-cascade citations here if desired; I have left this generic
because I do not have a high-confidence ADS key in hand for the most
appropriate references.}

A more powerful and increasingly mature class of statistics for
non-Gaussian fields is the wavelet scattering transform
\citep{CITATION_NEEDED_TBD}. Wavelet scattering coefficients are
translation-equivariant, rotationally invariant, multi-scale
descriptors that capture the morphological content of a field at a
hierarchy of spatial scales and orientation interactions. They are
strictly more informative than two-point statistics for fields with
non-Gaussian structure and have been used as field-level summary
statistics in ISM, CMB, and galaxy-morphology inference contexts
\citep{CITATION_NEEDED_TBD}. \CGK{Suggested references to fill in:
Allys+19, Cheng+20, R\'egaldo-Saint Blancard+20/+23; please verify
ADS keys and substitute.}

In the long term, the goal is field-level inference: training a neural
posterior estimator directly on simulated maps, with wavelet scattering
coefficients (or learned summary statistics) serving as the
intermediate representation. This program collapses the distinction
between forward modeling and inference: each \TIGRESS\ snapshot
produces, via the pipeline of \autoref{sec:disc_synth}, the same
observables that \PHANGS\ measures, and the neural estimator learns the
posterior over $\thetaenv$ (and eventually $\thetaphys$) directly from
the resolved map. The companion paper takes the first step toward this
goal using aperture-level summary statistics; the present suite is sized
and structured so that the same simulations can be reused as the
training set as field-level summaries are added.

\subsection{Connection to the Companion Inference Paper}
\label{sec:disc_paperII}

The simulation suite presented here is designed from the outset to
serve as the numerical prior and simulator for the simulation-based
inference (\SBI) framework developed in the companion paper
\CGK{(\citealt{CITATION_NEEDED_TBD}; Paper~II, in preparation)}. The
sampling design of \autoref{sec:sampling}, mapped through the KDE of
the \PHANGS\ joint distribution, defines the implicit prior $p(\thetaenv)$
under which the inference is performed. Summary statistics extracted
from each \TIGRESS\ snapshot---$\Ssfr$, $\fmol$, $\PDE$, $\seff$, and
the decomposed feedback yields $\Yth$ and $\Ynonth$---constitute the
simulated observables that are matched against \PHANGS\ aperture
measurements. The progressive extensibility of the Sobol design ensures
that the training set can be enlarged without invalidating any inference
performed on earlier subsets, and the reservation of $\fmol$ and
$\Ssfr$ as validation quantities means that the inferred posteriors can
be cross-checked against quantities that did not enter the design. The
present paper completes the construction of the simulator; the
companion paper applies it.


\begin{acknowledgments}
The authors acknowledge the use of Claude Code (Anthropic) as an AI-assisted
editing tool during the preparation of this manuscript. All scientific ideas,
physical interpretations, mathematical formulations, and intellectual content
are original contributions of the authors.
\end{acknowledgments}

\bibliography{../shared/references}

\end{document}
